Lidar point cloud technology has garnered significant attention due to its extensive applications in high-precision 3D reconstruction, particularly in fields such as autonomous driving, robot navigation, and cultural heritage conservation. With the advancement of point cloud data acquisition techniques, how to efficiently and accurately reconstruct 3D models from point cloud data has become a research hotspot. This paper investigates high-precision 3D reconstruction technology based on lidar point clouds, aiming to explore the performance differences between traditional algorithms and emerging implicit reconstruction methods, as well as their practical applications.

Firstly, the paper analyzes the basic theory and algorithm framework of point cloud 3D reconstruction, delving into traditional algorithms such as Alpha Shapes and Ball Pivoting Algorithm (BPA). By constructing a 3D display surface reconstruction algorithm, 3D surfaces are reconstructed from various open-source 3D point cloud datasets. The performance of different traditional algorithms is compared using multiple quantitative indicators (chamfer distance, Hausdorff distance), verifying the effectiveness of traditional methods in 3D point cloud data.

Secondly, the paper focuses on the implicit reconstruction method StEik based on neural signed distance functions and implements a 3D implicit surface reconstruction algorithm. Experimental results show that the StEik algorithm significantly enhances reconstruction accuracy and efficiency when dealing with complex point cloud data, especially in terms of noise robustness and detail reproduction, outperforming traditional methods.

To further verify the practical application capabilities of the algorithms, 3D printing technology is employed to create various objects, and precise point cloud datasets are acquired using a 3D scanner. By comparing the performance of different algorithms on the self-acquired datasets, it is found that the StEik algorithm demonstrates higher stability and adaptability in complex scenarios, providing reliable technical support for practical engineering applications.

In summary, this paper investigates 3D reconstruction techniques using both traditional and implicit methods. Through theoretical analysis and experimental validation, it provides research insights and technical pathways for high-precision 3D reconstruction, and holds significant importance in promoting the practical application of LiDAR point cloud reconstruction technology